% Options for packages loaded elsewhere
\PassOptionsToPackage{unicode}{hyperref}
\PassOptionsToPackage{hyphens}{url}
\PassOptionsToPackage{dvipsnames,svgnames,x11names}{xcolor}
\documentclass[
]{llncs}
\usepackage{xcolor}
\usepackage{amsmath,amssymb}
\setcounter{secnumdepth}{-\maxdimen} % remove section numbering
\usepackage{iftex}
\ifPDFTeX
  \usepackage[T1]{fontenc}
  \usepackage[utf8]{inputenc}
  \usepackage{textcomp} % provide euro and other symbols
\else % if luatex or xetex
  \usepackage{unicode-math} % this also loads fontspec
  \defaultfontfeatures{Scale=MatchLowercase}
  \defaultfontfeatures[\rmfamily]{Ligatures=TeX,Scale=1}
\fi
\usepackage{lmodern}
\ifPDFTeX\else
  % xetex/luatex font selection
  \setmainfont[]{Cambria}
\fi
% Use upquote if available, for straight quotes in verbatim environments
\IfFileExists{upquote.sty}{\usepackage{upquote}}{}
\IfFileExists{microtype.sty}{% use microtype if available
  \usepackage[]{microtype}
  \UseMicrotypeSet[protrusion]{basicmath} % disable protrusion for tt fonts
}{}
\makeatletter
\@ifundefined{KOMAClassName}{% if non-KOMA class
  \IfFileExists{parskip.sty}{%
    \usepackage{parskip}
  }{% else
    \setlength{\parindent}{0pt}
    \setlength{\parskip}{6pt plus 2pt minus 1pt}}
}{% if KOMA class
  \KOMAoptions{parskip=half}}
\makeatother
\setlength{\emergencystretch}{3em} % prevent overfull lines
\providecommand{\tightlist}{%
  \setlength{\itemsep}{0pt}\setlength{\parskip}{0pt}}
\usepackage{bookmark}
\IfFileExists{xurl.sty}{\usepackage{xurl}}{} % add URL line breaks if available
\urlstyle{same}
\hypersetup{
  pdftitle={SoK: The Shelf Life of Privacy Guarantees},
  pdfauthor={Anonymous Submission to FC 2027},
  colorlinks=true,
  linkcolor={black},
  filecolor={Maroon},
  citecolor={Blue},
  urlcolor={black},
  pdfcreator={LaTeX via pandoc}}

\title{SoK: The Shelf Life of Privacy Guarantees}
\author{Anonymous Submission to FC 2027}
\institute{Anonymized for review}
\date{}

\begin{document}
\maketitle

\subsection{Time-Indexed Reconstruction Bounds over Fixed
Archives}\label{time-indexed-reconstruction-bounds-over-fixed-archives}

\subsection{Abstract}\label{abstract}

Privacy guarantees over recorded data degrade in a specific and
under-formalised way: the record is fixed at disclosure time, while the
adversary's decoding capability keeps improving. At least four
literatures observe this independently, as harvest-now-decrypt-later, as
retroactive de-anonymisation, as cryptoperiod practice, and as the
motivation for everlasting-privacy constructions, but no common formal
object connects them. This paper systematises these literatures under a
time-indexed reconstruction ratio R(t) = (C\_S(t) + C\_M(t))/H(X), the
ratio of the effective capacities of the channels leaving a two-agent
disclosure architecture, evaluated against the strongest adversary class
available at time t, to the entropy of the protected source. Under two
named preconditions, non-collusion between the channels and evaluation
against a declared adversary class, the capacity sum bounds the joint
leakage and a Fano-type error floor holds; the strict bound R(t)
\textless{} 1 holds exactly when a third, deliberately separate
condition holds, the declared capacity-deficit condition C\_S(t) +
C\_M(t) \textless{} H(X), and it is this deficit condition, not the
architecture, that expires at the shelf life t* = sup\{t : R(t)
\textless{} 1\}. The drift itself decomposes into two mechanisms kept
distinct throughout: method-and-compute growth, which erodes
computationally protected channels and moves certified estimates but
cannot raise leakage through an information-theoretic bound, and
side-information growth, which erodes even information-theoretic floors
through the residual entropy of the source given the adversary's
accumulating background. We survey the families the ratio organises,
including two verified 2026 instances of the archival regime, state the
adversary-ordering formalisation this systematisation requires as an
explicit programme, state the existence-leak conjecture for capability
attestations in its operational form, that conditioning on attested
feasibility provably reduces the adversary's search cost (its bare
inequality I(feasibility; method) \textgreater{} 0 is shown vacuously
true under the naive prior model), with the 2026 withheld-capability
episode as its single verified instance, and report one countermeasure
direction as an unvalidated open problem. Conjectures are stated as
conjectures with proof obligations; the survey's unification claim is
made to our knowledge and its known referee risks are stated in the
threats section.

\subsection{1 Introduction}\label{introduction}

A confidentiality guarantee issued today over data recorded today is,
implicitly, a claim about every future adversary who will ever hold the
recording. The wiretap lineage that grounds information-theoretic
secrecy fixes the adversary together with the channel (Wyner 1975;
Leung-Yan-Cheong and Hellman 1978; Csiszár and Körner 1978) and
therefore cannot express the deployed-systems situation in which the
channel output is archived and the adversary class improves against it
year over year. The situation itself is well documented. The
harvest-now-decrypt-later literature treats it as an economic threat
model (Blanco-Romero et al.~2026). The retroactive de-anonymisation
record shows datasets released as safe becoming unsafe when auxiliary
data and methods arrived later (Narayanan and Shmatikov 2008; Gymrek et
al.~2013; Erlich et al.~2018). The consent-side statement of the same
phenomenon is the end run: inference from accumulating background data
circumvents anonymity and consent protections alike, so a guarantee's
effective strength is a function of the evolving external data
environment rather than of its issuing assumptions (Barocas and
Nissenbaum 2014). Key-management practice bounds key lifetimes
administratively (Barker 2020). The everlasting-privacy line constructs
protocols whose privacy survives later, computationally unbounded
adversaries, precisely because its authors did not trust computational
assumptions to age (Moran and Naor 2006; Haines et al.~2023). These
strands share a phenomenon and do not share a formalism.

This paper's position is that the shared object is a ratio. Let H(X) be
the entropy of the protected source, fixed at disclosure time, and let
C\_S(t) and C\_M(t) be the effective capacities of the two channels
leaving a disclosure architecture that splits data flow between a
boundary agent S and a delegation agent M, evaluated against the
strongest adversary class available at time t. The reconstruction ratio
is R(t) = (C\_S(t) + C\_M(t))/H(X), and the shelf life of the guarantee
is t* = sup\{t : R(t) \textless{} 1\}. The formal content of this paper
is a decomposition of what makes R(t) \textless{} 1 a guarantee at all:
two preconditions (non-collusion and a declared adversary class) license
the capacity-sum bound on joint leakage and a Fano-type error floor, and
a separate, declarable, measurable capacity-deficit condition licenses
the strict bound. The preconditions do not by themselves place the ratio
below one. What expires at t* is the deficit condition, not the
architecture.

We write this as a systematisation of knowledge rather than a results
paper for a stated reason: the value of the ratio is that five
literatures that do not cite one another become instances of one object,
and the honest status of the object's own apparatus is mixed. The bound
decomposition is conditional and elementary; the adversary ordering it
presupposes is a formalisation programme, not a finished definition; the
drift of R(t) and the leakage of capability attestations are conjectures
carried in the research programme's conjecture register with stated
evidence status; and the countermeasure direction we can name has no
measurement behind it at all. Each of these appears below at exactly
that status. A result against the model's prediction, should one arrive,
is filed to the register with the same prominence as a confirmation, and
this paper commits to that in writing.

\subsubsection{1.1 Contributions}\label{contributions}

\begin{enumerate}
\def\labelenumi{\arabic{enumi}.}
\tightlist
\item
  \textbf{A derived shelf life (Section 3).} The definition of the
  reconstruction ratio R(t) = (C\_S(t) + C\_M(t))/H(X), capacities
  evaluated against the strongest adversary class available at time t,
  and of the shelf life t* = sup\{t : R(t) \textless{} 1\} as a quantity
  derived from an architecture and an adversary class, where prior art
  states lifetimes as declared requirements (Mosca 2018; Barker 2020) or
  compares declared lifetimes against declared horizons.
\item
  \textbf{The deficit decomposition (Section 3).} A separation of
  obligations within the guarantee: the two preconditions, non-collusion
  and a declared adversary class, are attestable properties of a
  deployment and license the capacity-sum bound on joint leakage and a
  Fano-type error floor, while the capacity-deficit condition licensing
  the strict bound R \textless{} 1 is a numerical comparison that no
  architectural property entails, and it is the deficit condition, not
  the architecture, that expires. The wiretap lineage owns the floor
  machinery; the novelty claim is confined to isolating a declared,
  measurable deficit condition as the time-indexed part of a privacy
  guarantee, which to our knowledge no prior work does.
\item
  \textbf{An adversary-ordering programme (Section 4).} The
  formalisation programme for ``adversary class at time t'': an ordered
  family \{D\_t\} with D\_t ⊆ D\_t' for t ≤ t', under which the
  monotonicity of R(t) is structural (Lemma 1) and only the rate is
  empirical. The family is grounded at definition level in the
  gain-function framework of quantitative information flow, where
  adversary models are currently per-scenario and unordered in time, and
  the programme remains contingent: its residual obligations are named
  exactly in Section 4.1, and if they fail the monotonicity demotes to a
  modelling assumption.
\item
  \textbf{The unification (Section 5).} To our knowledge, no prior
  systematisation organises harvest-now-decrypt-later, retroactive
  de-anonymisation, quantum-horizon estimation, cryptoperiod practice,
  continual-observation differential privacy, and everlasting-guarantee
  constructions as instances of a single time-indexed reconstruction
  bound; the nearest prior SoK is confined to e-voting (Haines et
  al.~2023) and the nearest survey of long-term guarantees covers
  integrity, not confidentiality (Vigil et al.~2015).
\item
  \textbf{An existence-leak law (Section 6).} The conjecture that a
  zero-knowledge proof of capability feasibility leaks an upper bound on
  reconstruction or search difficulty, scoped to capability
  attestations, floored by the leakage-resilient zero-knowledge
  impossibility of Garg, Jain and Sahai (2011), and instantiated by a
  verified 2026 episode. Its displayed inequality, I(feasibility;
  method) \textgreater{} 0, is shown in Section 6.1 to be vacuously true
  under the naive prior model; the contribution is the conjecture's
  operational form and its proof obligation, an attestation class in
  which conditioning on attested feasibility provably and quantifiably
  reduces the adversary's search cost.
\item
  \textbf{An attestation discount (Section 6).} The conjecture that
  public feasibility attestations shorten the effective migration
  horizon of Mosca-style planning independently of any actual attack,
  Z\_b' = Z\_b - D(a).
\item
  \textbf{A named open problem (Section 7).} The divergence
  countermeasure, reconstruction error growing as e\^{}(lambda t) along
  a diverging data trajectory, presented strictly as an open problem:
  lambda is unmeasured, and we found neither anticipation nor support in
  the literature.
\end{enumerate}

\subsection{2 Systematization Method}\label{systematization-method}

\textbf{Corpus and verification.} The survey beneath this paper swept
nine categories: harvest-now-decrypt-later; migration-timeline analysis
in the Mosca lineage; cryptoperiod and key-management practice;
retroactive de-anonymisation; quantum resource estimation for
elliptic-curve discrete logarithms; the wiretap lineage's treatment of
time; everlasting and long-term guarantees, including prior SoKs on
expiring guarantees; differential privacy under continual observation;
and adjacent work on existence leaks, leakage-resilient zero knowledge,
and quantitative information flow. Forty-one works and primary records
were resolved to a live publisher page, index entry, or DOI, including
the nearest published neighbour to this paper's central definition,
which is engaged in Section 5.1 (Kagai, Branch, But and Allen 2025). One
further candidate work could not be bibliographically resolved (no
resolvable venue or author list); it is excluded from citation and
flagged where relevant (Sections 5.1 and 8).

\textbf{Inclusion criterion.} A work enters the survey if it gives
confidentiality, anonymity, or reconstruction a time axis: a guarantee
that decays, a lifetime that is declared or compared, an adversary that
improves, or a construction whose purpose is to survive such
improvement.

\textbf{Organising question.} For each family we ask where time enters
the guarantee. Three answers occur: time enters through the releases
(more output accumulates against a fixed adversary, the
continual-observation regime); through the channel (the channel itself
varies while the adversary class is fixed, the fading regime); or
through the adversary (the recorded output is fixed and the decoder
class grows, the archival regime). The ratio R(t) formalises the third
answer, which is the one the surveyed families instantiate and the one
the formal lineages have not treated as an object (Section 5.7).

\textbf{Claim discipline.} Every claim in this paper is fenced by a
prior-art novelty ledger built before drafting: assertions presented as
new are stated no wider than the ledger's wording; assertions the
literature anticipates are framed as synthesis with the anticipating
work cited. Conjectures are stated as conjectures with their register
identifier in the research programme's conjecture register and a
qualitative statement of their evidence status; no conjecture is
presented as a result. Negative findings, including the absence of
support for the open problem of Section 7, are reported with the same
prominence as positive ones.

\subsection{3 Model and Preliminaries: the Time-Indexed Reconstruction
Bound}\label{model-and-preliminaries-the-time-indexed-reconstruction-bound}

\subsubsection{3.1 Notation and system
model}\label{notation-and-system-model}

Random variables are uppercase; H(X) is Shannon entropy in bits and
I(X;Y) mutual information, with standard definitions per Cover and
Thomas (2006). The source X is the data subject's protected information,
taking values in a finite alphabet 𝒳 with \textbar 𝒳\textbar{} ≥ 3 and
H(X) \textgreater{} 0, and with H(X) fixed by the source at disclosure
time. The disclosure architecture routes data flow through two agents: a
boundary agent S, which mediates the subject's outward disclosures, and
a delegation agent M, which performs delegated processing.\footnote{The
  disclosure architecture and its running model are maintained publicly
  at https://agentprivacy.ai/model. This is the only non-literature URL
  this paper carries, per its claim discipline.} Their outward-facing
outputs are Y\_S and Y\_M. An adversary observes the outputs available
to its class and attempts to reconstruct X; for an estimator X̂ of X from
the observed outputs, P\_e := Pr{[}X̂ ≠ X{]} denotes its error
probability. C\_S(t) and C\_M(t) denote the effective capacities of the
two channels with respect to the strongest adversary class available at
time t; capacity here is the adversary-relative effective capacity of
each output channel for information about X, a notion whose
formalisation is taken up in Section 4 (Definition 3).

\subsubsection{3.2 Assumptions}\label{assumptions}

\textbf{Assumption 1 (non-collusion).} I(Y\_S; Y\_M \textbar{} X) = 0,
and no third channel carries the inter-agent residue: the two outputs
are conditionally independent given the source, and the architecture
provides no side channel that recreates their joint structure.

\textbf{Assumption 2 (declared adversary class).} The capacities C\_S
and C\_M are evaluated against a stated adversary class, declared
together with any bound derived from them; against that class the
channel leakages about the source are bounded as I(X; Y\_S) ≤ C\_S and
I(X; Y\_M) ≤ C\_M. (Definition 3 in Section 4 states what ``evaluated
against a class'' is required to mean for these bounds to be well
defined; no result in this section uses more than the two displayed
inequalities.)

\subsubsection{3.3 The conditional floor and the deficit
condition}\label{the-conditional-floor-and-the-deficit-condition}

\textbf{Proposition 1 (capacity-sum bound and conditional error floor).}
Let X take values in a finite alphabet 𝒳 with \textbar 𝒳\textbar{} ≥ 3,
and let Assumptions 1 and 2 hold, so that I(X; Y\_S) ≤ C\_S and I(X;
Y\_M) ≤ C\_M against the declared class. Write R\_max := (C\_S +
C\_M)/H(X). Then:

\begin{enumerate}
\def\labelenumi{(\alph{enumi})}
\item
  \emph{(capacity-sum bound)} I(X; Y\_S, Y\_M) ≤ I(X; Y\_S) + I(X; Y\_M)
  ≤ C\_S + C\_M, with equality in the first inequality if and only if
  I(Y\_S; Y\_M) = 0;
\item
  \emph{(error floor)} every estimator X̂(Y\_S, Y\_M) satisfies
\end{enumerate}

\begin{quote}
P\_e ≥ (H(X) − C\_S − C\_M − 1) / log(\textbar 𝒳\textbar{} − 1),
\end{quote}

and if moreover H(X) ≥ log(\textbar 𝒳\textbar{} − 1) (in particular if X
is uniform on 𝒳), then

\begin{quote}
P\_e ≥ 1 − R\_max − 1/H(X).
\end{quote}

\emph{Proof.} (a) By the chain rule, I(X; Y\_S, Y\_M) = I(X; Y\_S) +
I(X; Y\_M \textbar{} Y\_S). The interaction identity I(X; Y\_M) − I(X;
Y\_M \textbar{} Y\_S) = I(Y\_S; Y\_M) − I(Y\_S; Y\_M \textbar{} X) (both
sides equal the co-information of the triple) gives, under Assumption 1,
I(X; Y\_M \textbar{} Y\_S) = I(X; Y\_M) − I(Y\_S; Y\_M) ≤ I(X; Y\_M),
with equality exactly when I(Y\_S; Y\_M) = 0; Assumption 2 supplies the
outer bound. (b) Fano's inequality (Fano 1961; Cover and Thomas 2006)
gives H(X \textbar{} Y\_S, Y\_M) ≤ h(P\_e) + P\_e
log(\textbar 𝒳\textbar{} − 1) ≤ 1 + P\_e log(\textbar 𝒳\textbar{} − 1),
with h the binary entropy. Substituting H(X \textbar{} Y\_S, Y\_M) =
H(X) − I(X; Y\_S, Y\_M) ≥ H(X) − C\_S − C\_M by (a) and rearranging
gives the first display. For the second: when H(X) − C\_S − C\_M − 1 ≥ 0
and H(X) ≥ log(\textbar 𝒳\textbar{} − 1), the first display's right side
is at least (H(X) − C\_S − C\_M − 1)/H(X) = 1 − R\_max − 1/H(X); when
H(X) − C\_S − C\_M − 1 \textless{} 0 the second display's right side is
negative and the bound holds trivially. ∎

\emph{Attribution.} Neither part of Proposition 1 is new, and this paper
claims neither as a contribution. Part (a) is the standard
sub-additivity of leakage under conditional independence, the summation
regime of the wiretap lineage (Wyner 1975; Leung-Yan-Cheong and Hellman
1978), whose general single-eavesdropper form is the broadcast channel
with confidential messages (Csiszár and Körner 1978). The failure of the
summed bound under collusion needs no citation: it is the elementary
fact that joint information can exceed the sum of marginal informations
once the outputs correlate beyond what the source explains, and the
interaction identity in the proof of (a) locates the excess exactly
(Section 8, assumption failure). Part (b) is the Fano converse in its
source-coding form (Fano 1961; Cover and Thomas 2006). The proposition
is recorded as a restated known-bound pair with its exact hypotheses;
the claim this paper makes is confined to the decomposition reading
below. A precision note: the floor is often quoted in the rounded form
P\_e ≥ 1 − R\_max; the exact converse carries the additive 1/H(X) term
and the alphabet-entropy hypothesis of (b), and only the exact form is
asserted here.

The decomposition that this paper isolates as its second contribution is
what Proposition 1 does not say. The two preconditions do not by
themselves place R\_max below one. The strict bound R\_max \textless{} 1
holds exactly when, additionally, the \textbf{capacity-deficit
condition} holds:

\begin{quote}
C\_S + C\_M \textless{} H(X),
\end{quote}

a measurable, declarable, numerical fact about a given system and
adversary class, not a consequence of the architecture. (At the level of
hypotheses the equivalence is definitional, since R\_max is the deficit
ratio; its content is the separation of obligations: Assumptions 1 and 2
are attestable properties of a deployment, while the deficit is a
numerical comparison that no architectural property entails.) The
deficit condition is deliberately not listed as a third assumption: two
conditionally independent channels of sufficient combined capacity
satisfy both assumptions with R\_max at or above one, and at R\_max ≥ 1
the error floor of Proposition 1(b) is vacuous, which is exactly why the
deficit condition must be declared alongside the preconditions wherever
the strict bound is claimed. Per this discipline, no passage of this
paper states R \textless{} 1 without its preconditions and its time
index.

\subsubsection{3.4 The time index and the shelf
life}\label{the-time-index-and-the-shelf-life}

\textbf{Definition 2 (reconstruction ratio and shelf life).} Fix a
disclosure time t\_0 and, for each t ≥ t\_0, let C\_S(t) and C\_M(t) be
the effective capacities of the two channels against the strongest
adversary class available at time t, with H(X) fixed by the source at
t\_0. The reconstruction ratio is

\begin{quote}
R(t) = (C\_S(t) + C\_M(t))/H(X), t ≥ t\_0,
\end{quote}

and the shelf life of the guarantee is

\begin{quote}
t* = sup\{t ≥ t\_0 : R(t) \textless{} 1\},
\end{quote}

with two conventions: t* = ∞ if R(t) \textless{} 1 for all t ≥ t\_0, and
t* undefined (the guarantee never held) if \{t ≥ t\_0 : R(t) \textless{}
1\} is empty.

\textbf{Remark 1 (what the supremum does and does not assert).}
Definition 2 takes the supremum reading deliberately, and it is the
weaker of the two available readings. What it asserts is exact: by
definition of the supremum, R(t) ≥ 1 for every t \textgreater{} t\emph{,
so no time after t} satisfies the deficit condition. What it does not
assert is equally exact: absent further hypotheses, R(t) \textless{} 1
need not hold for all t ∈ {[}t\_0, t*{]}, since a non-monotone R may
reach one and return below it before t\emph{; the interval up to t} is
not certified, only the region beyond t* is excluded. The first-crossing
quantity inf\{t ≥ t\_0 : R(t) ≥ 1\} is not used anywhere in this paper,
and no statement below depends on the two readings coinciding. They
coincide exactly when \{t : R(t) \textless{} 1\} is an interval with
left endpoint t\_0, in particular whenever R is non-decreasing; and
whether R is non-decreasing is precisely the ordering question of
Section 4, where Lemma 1 derives monotonicity from the ordered-family
hypothesis, which is a hypothesis about adversary classes, not an
established fact. The behaviour of t* under a non-monotone R is
therefore an open question of the formalisation programme, recorded here
rather than assumed away.

\textbf{Remark 2 (availability semantics).} Definition 2's phrase ``the
strongest adversary class available at time t'' admits two readings, and
this paper fixes the first: D\_t is the class of decoders realisable by
some adversary at time t, whether or not the enabling methods are
published; the publicly documented capability trajectory is an auditor's
lower estimate of D\_t. The archival regime forces this reading: an
adversary holding the recording decodes it with whatever capability it
has, disclosed or not, so a ratio evaluated against publicly available
capability alone would certify a guarantee that an undisclosed
capability had already voided. Two consequences follow and are used
below. First, R(t) and t* are properties of realisable capability and
are therefore not directly observable: any computed shelf life is an
estimate of t* built from the auditor's lower estimate of \{D\_t\}, and
the estimation error is in the unsafe direction, since understating D\_t
overstates t*. Second, a disclosure or attestation of an
already-realised capability moves the auditor's estimate of R(t), not
R(t) itself; where publication additionally extends what other
adversaries can realise, that propagation is a distinct, subsequent
growth of D\_t. Conjecture 1's coupling of drift to release events is a
claim about this propagation mechanism; Section 5.8's
withheld-capability episode is read against these semantics in place,
and the attestation discount of Conjecture 4 acts on the planner's
estimate of the horizon, not on the derived crossing.

\textbf{Remark 2b (two mechanisms of growth, and two clocks).} The
family \{D\_t\} grows through two mechanisms that this survey keeps
distinct, because they erode different things. The first is
method-and-compute growth: better estimators, better solvers, more
resources. For computationally protected channels (the cryptoperiod and
harvest-now-decrypt-later families of Section 5.3) this mechanism
genuinely erodes the protection; for a channel whose leakage bound is
information-theoretic it cannot, since such a bound already holds
against unbounded computation, and what it moves instead is the
certified estimate of the capacities, the auditor's quantity of Remark
2. The second is side-information growth: the adversary's auxiliary
corpus and priors accumulate, so every decoder in the class works from
more background than before. This mechanism erodes even
information-theoretic floors, and it is the one the retroactive
de-anonymisation record of Section 5.4 instantiates. The companion
theory paper carries the second mechanism as proved material: with B\_t
the adversary's accumulating background (accumulation ordered in t,
Markov B\_t → X → T against source X and transcript T), the disclosed
transcript's leakage bound survives conditioning while the Fano floor
erodes through the residual entropy H(X \textbar{} B\_t), so the
informed deficit condition reads C\_S + C\_M \textless{} H(X \textbar{}
B\_t) and the erosion requires no growth in what the transcript reveals.
Where a passage below speaks of drift without qualifying its mechanism,
the load-bearing mechanism for information-theoretically bounded
disclosures in the archival regime is the second.

The decomposition of Section 3.3 makes the time-dependence precise: the
deficit condition, not the preconditions, is the time-indexed quantity.
R(t) can cross one with both preconditions intact; what expires at t* is
the deficit condition, not the architecture. This is the formal shape of
the archival regime: the recording is fixed, the assumptions under which
it was made can remain true forever, and yet no time beyond t* satisfies
the declared deficit against the then-strongest adversary class.

The distinction against declared lifetimes is the content of the first
contribution. A cryptoperiod (Barker 2020) and the security shelf life x
of Mosca's planning inequality (Mosca 2018) are requirements stated by
the data owner; the vulnerability-window comparisons of the
harvest-now-decrypt-later literature compare a declared confidentiality
lifetime against a projected decryption horizon. In Definition 2, t* is
derived from the architecture's declared capacities and the
adversary-class trajectory: the inputs are declared quantities (Section
8), and the crossing time, given them, is computed, not chosen.

\subsection{4 The Ordered Adversary Family: a Formalisation
Programme}\label{the-ordered-adversary-family-a-formalisation-programme}

\subsubsection{4.1 The programme}\label{the-programme}

Definition 2 uses the phrase ``strongest adversary class available at
time t''. For the ratio to be an object rather than a metaphor, that
phrase requires formalisation as an ordered family of decoder classes
\{D\_t\}, with D\_t ⊆ D\_t' whenever t ≤ t', so that later adversary
classes can do everything earlier ones could. The natural grounding is
quantitative information flow, where adversaries are modelled by gain
functions and leakage is quantified against them, with capacity as a
worst case over gain functions (Alvim et al.~2012; Alvim et al.~2020).
We give the grounding at definition level, prove the one structural
lemma it licenses, and then state exactly what remains open; the
boundary between the structural and the empirical content below is the
boundary of the claim.

\textbf{Definition 3 (decoder class; effective capacity relative to a
class).} A \emph{decoder class} D is a set of adversary models in the
sense of quantitative information flow: each element d ∈ D specifies a
gain function g\_d (Alvim et al.~2012) together with the strategy
resources permitted in realising a g\_d-optimal attack. A
\emph{capability-indexed family} is a map t ↦ D\_t from times to decoder
classes. Let ℓ be a leakage functional, valued in bits, defined per
adversary model, and subject to the \emph{compatibility requirement}:
for every model d and every estimator X̂ realisable by d from a channel
output Y, I(X; X̂(Y)) ≤ ℓ\_d(X → Y). The \emph{effective capacity} of a
channel with output Y against the class D is

\begin{quote}
C(Y; D) = sup\_\{d ∈ D\} ℓ\_d(X → Y),
\end{quote}

and Definition 2's capacities are C\_S(t) = C(Y\_S; D\_t), C\_M(t) =
C(Y\_M; D\_t). The compatibility requirement is the minimum a witness
for the functional ℓ must satisfy, and what it secures is exactly
per-channel: every estimator realisable within the class from a single
channel output Y extracts at most C(Y; D) bits about X. Two things it
does not secure, and this definition does not claim, are Assumption 2's
displayed inequalities read with these effective capacities, which
remain a substantive assumption on the class-channel pair and can fail
precisely when the effective capacity sits below the true mutual
information, and the extension of Proposition 1's floor to estimators of
X from the joint output (Y\_S, Y\_M), which needs in addition a
sub-additivity clause for ℓ across conditionally independent channels
that is nowhere required here. What a finished definition must
additionally deliver, per item (ii) of the empirical residue below, is
that joint-estimator floor: until a functional delivering it is
exhibited, Proposition 1's floor evaluated with effective capacities is
a desideratum of the programme, not a consequence of the compatibility
requirement.

\textbf{Lemma 1 (structural monotonicity).} Let t ↦ D\_t be a
capability-indexed family satisfying the \emph{ordering hypothesis}:
D\_t ⊆ D\_t' whenever t ≤ t'. Then C(Y; D\_t) ≤ C(Y; D\_t') for each
channel, hence R(t) is non-decreasing; consequently \{t ≥ t\_0 : R(t)
\textless{} 1\} is, when nonempty, an interval with left endpoint t\_0;
and whenever both \{t ≥ t\_0 : R(t) \textless{} 1\} and \{t ≥ t\_0 :
R(t) ≥ 1\} are nonempty, so that t* is defined and finite, the supremum
of Definition 2 equals the first-crossing quantity inf\{t ≥ t\_0 : R(t)
≥ 1\}. (If only the sub-unity set is nonempty, t* = ∞; if the sub-unity
set is empty, t* is undefined by Definition 2's convention while the
first-crossing quantity is t\_0, so the equality clause is not asserted
there.)

\emph{Proof.} A supremum over a larger set is no smaller, so each
effective capacity is non-decreasing in t under the ordering hypothesis,
and R(t), a sum of non-decreasing functions over the positive constant
H(X), is non-decreasing. For non-decreasing R, if R(s) \textless{} 1 and
t\_0 ≤ t ≤ s then R(t) \textless{} 1, so the sub-unity set is an
interval with left endpoint t\_0. Let τ = sup\{t : R(t) \textless{} 1\}
and σ = inf\{t : R(t) ≥ 1\}, both sets nonempty. For every t
\textgreater{} τ, R(t) ≥ 1, so σ ≤ τ + δ for every δ \textgreater{} 0,
giving σ ≤ τ. Conversely, for every t \textless{} τ there is, by
definition of the supremum, an s ∈ (t, τ{]} with R(s) \textless{} 1,
hence R(t) ≤ R(s) \textless{} 1 by monotonicity; so every member of \{t
: R(t) ≥ 1\} is at least τ, giving σ ≥ τ. ∎

Under the ordering hypothesis, then, the monotone non-decrease of R(t)
against a fixed archive is structural, the two readings of t* in Remark
1 coincide, and the empirical content of any drift claim reduces to the
rate of increase, not its sign.

\textbf{The empirical residue, stated exactly.} What Lemma 1 does not
supply is the programme, and it is threefold. (i) \emph{The ordering
hypothesis itself.} That later adversary classes can do everything
earlier ones could is plausible under non-regression of published
methods, but it is a modelling assumption about capability, not a
theorem; capability can in principle regress (methods lost, compute
repriced), and no result above survives such regress. (ii) \emph{The
functional ℓ.} The compatibility requirement of Definition 3 is stated,
not discharged: the QIF candidates (g-leakage, Bayes capacity and its
logarithm) quantify leakage against gain models, but a bit-valued
functional calibrated against source entropy, uniformly over a class,
such that the Fano argument of Proposition 1 survives the substitution
of effective for information-theoretic capacities against estimators
from the joint output, for instance via a sub-additivity clause across
conditionally independent channels, has not been exhibited; this is the
exact missing step between Definition 3 and a finished definition, and
it is the delivery named at the definition's close. (iii)
\emph{Instantiation.} No published adversary series is presently
expressed as a nested gain-function family: QIF adversary models are
per-scenario and unordered in time, continual-observation differential
privacy fixes the adversary while releases grow (Dwork et al.~2010), and
the quantum resource-estimate literature provides point estimates
without treating the series as an ordered process (Sections 5.7 and
5.8). The programme is contingent in the following sense: if (i) fails
for the adversary classes of interest, or (ii) admits no witness, then
decoder classes cannot be coherently ordered as gain-function families,
the monotonicity of Lemma 1 demotes from a structural fact to a
modelling assumption, this paper's fourth section demotes with it to a
discussion, and Remark 1's supremum reading becomes load-bearing rather
than merely cautious.

\subsubsection{4.2 The drift conjecture}\label{the-drift-conjecture}

That guarantees erode against fixed archives is not a conjecture; it is
the shared observation of the harvest-now-decrypt-later corpus, the
everlasting-privacy motivation, and the retroactive de-anonymisation
record surveyed in Section 5. What the research programme adds, and
holds at conjecture status, is a claim about the driving process:

\textbf{Conjecture 1 (informational-capability drift; programme register
C82 as re-typed 2026-07-17, moderate confidence, rate unestablished).}
Adversary informational capability grows against fixed archives: the
linkage corpus and side priors accumulate along calendar time, shrinking
H(X \textbar{} B\_t) while nothing is added to the archive and no action
of the data subject is involved; R(t) drifts upward on a schedule, and
every static reconstruction guarantee has a finite shelf life t\emph{.
Frontier-model releases enter only informationally, through improved
extraction of linkage from existing corpora, never as computation
against the information-theoretic guarantee, which is compute-saturated.
In the terms of Definition 3, the claim is that the realised family
trajectory t ↦ D\_t grows as background accumulates and that the growth
strictly increases the effective capacities against fixed archives.
}Proof obligation:* the erosion form is proven conditional in the
companion (Definition 3.9 and Corollary 5.4b); what remains conjectural
is the rate: a formalised accumulation model and at least one measured
capability series showing capacity increase attributable to background
growth against a fixed archive; the two verified 2026 instances of
Section 5.8 are consistent with the conjecture but constitute a single
year's cluster, not a series. \emph{Mechanism typing (Remark 2b):} the
component of any such drift that bears on an information-theoretically
bounded disclosure is the informational one; the register's re-typing of
the conjecture to that mechanism was executed 2026-07-17.

\subsection{5 The Surveyed Families Under the
Bound}\label{the-surveyed-families-under-the-bound}

Each family below is read against Definition 2: which component of R(t)
does the family vary, declare, or construct against?

\subsubsection{5.1 Harvest-now-decrypt-later: the deficit condition as
economics}\label{harvest-now-decrypt-later-the-deficit-condition-as-economics}

The HNDL literature is the archival regime stated as a threat economics.
Blanco-Romero et al.~(2026) reframe HNDL as an economic problem and
quantify adversary storage and collection costs across TLS 1.2/1.3, QUIC
and SSH with an open testbed; the adversary is indexed by cost, not by
decoder class, and no shelf-life quantity is derived from architecture.
A Federal Reserve Board analysis of distributed-ledger data (FEDS
2025-093) notes that ledger immutability makes the archive term
irreducible, which in the terms of Section 5.6 is the statement that an
archive that cannot be erased has no t-independent term. In R(t) terms,
HNDL is the observation that C(t) will cross the deficit threshold for
recorded ciphertext and that collection is rational before the crossing;
what the family does not contain is a ratio against source entropy or a
derived t*.

The nearest published neighbour to Definition 2's purpose is the
temporal-risk model of Kagai, Branch, But and Allen (2025), which
presents itself as the first formal model of the HNDL adversary: a
resource-accumulating attacker specified by collection capability,
deferred decryption power and a temporal horizon, with compromise stated
as the lifetime-versus-horizon condition Ld \textgreater{} Ha, a
time-indexed compromise probability Pr\{Ha(t) ≥ Ld\} (written R(t)
there; the symbol names a probability, not this paper's
capacity-to-entropy ratio), break-time curves parameterised by hardware
trajectories (logical qubits, error-correction overhead), and a sectoral
exposure window W = max(0, Ld − Ha(t\_mig)) under which hybrid and
forward-secure mechanisms reduce the modelled risk horizon by over
two-thirds. Read against Definition 2, the work is the archival regime
formalised as risk arithmetic over declared quantities: the
confidentiality lifetime Ld is a regulatory or sectoral retention input,
the horizon Ha(t) is a projected scalar under scenario curves, and
neither is derived from the disclosing architecture; there is no ratio
against source entropy, no derived shelf life, and no ordering of
decoder classes, the adversary being indexed by a single horizon rather
than by a growing class of estimators. It therefore instantiates this
family rather than anticipating the object: t* is the computed
counterpart of its projected Ha, derived from declared channel
capacities and an adversary trajectory rather than adopted from forecast
baselines. A further work, a preprint reported to formalise the same
lifetime-versus-horizon condition, is bibliographically unresolved (no
resolvable venue or author list) and is therefore not cited; it is the
one open referee risk of this kind for the unification claim (Section
8).

\subsubsection{5.2 Migration timelines: the horizon as elicited
input}\label{migration-timelines-the-horizon-as-elicited-input}

Mosca's inequality, x + y \textgreater{} z (security shelf life plus
migration time exceeding the collapse horizon), is the planning rule of
the transition literature (Mosca 2018), and the Quantum Threat Timeline
series supplies z as expert-elicited probability distributions (Mosca
and Piani 2024), with the sharpest upward shift of the series in its
most recent edition (Mosca and Piani 2025). All three quantities are
declared or elicited; none is derived from system structure. In the
terms of this paper, Definition 2's t* plays the role of z, not of x: it
is the architecture's computed collapse horizon for a given adversary
trajectory. The cryptographic-agility literature systematised by Näther
et al.~(2024) addresses y, the cost of changing cryptography under
threat evolution; it measures switching cost, not any property of the
protected data trajectory, which distinguishes it from the open problem
of Section 7. Public capability rulers such as the graded ECDLP
challenge ladder of Dallaire-Demers, Doyle and Foo (2025) instrument the
adversary trajectory itself, which sits upstream of any particular
channel's effective capacity, and are, in this paper's terms,
measurement infrastructure for the adversary family of Section 4.

\subsubsection{5.3 Cryptoperiods: the declared
lifetime}\label{cryptoperiods-the-declared-lifetime}

NIST key-management practice bounds key lifetimes administratively,
typically at one to three years (Barker 2020). This is a weak
anticipation of the shelf-life idea: institutional practice already
treats protection as time-bounded. The distinction is categorical, not
quantitative: cryptoperiods bound key usage, and rotation does nothing
for already-exfiltrated ciphertext, which is exactly the regime R(t)
addresses; the guarantee over the archive is untouched by rotating the
key that produced it.

\subsubsection{5.4 Retroactive de-anonymisation: measured instances of
capability
drift}\label{retroactive-de-anonymisation-measured-instances-of-capability-drift}

This record is the empirical backbone of the archival regime. Netflix
Prize microdata released as perturbed and safe was re-identified when
auxiliary data arrived (Narayanan and Shmatikov 2008). Anonymised
genomes were re-identified through public genealogy databases that did
not exist at scale when the data was shared (Gymrek et al.~2013), and
long-range familial search made a majority of a large population
identifiable through relatives' enrolment, coverage growing without any
action of the subjects (Erlich et al.~2018), a sharp instance of
third-party growth eroding a guarantee whose issuing assumptions never
changed. Generative models estimate re-identification success against
incomplete releases (Rocher et al.~2019), and mobility traces identify
most individuals from four spatio-temporal points (de Montjoye et
al.~2013), which calibrates how small H(X) can be relative to auxiliary
channels. Read against Section 4: Rocher et al.~is the nearest
methodological neighbour to the decoder-class idea, predicted
reconstruction success under an explicit attacker model, but the
attacker model is static; the family \{D\_t\} is exactly what is
missing. Read against Remark 2b, this family is the record of the second
mechanism: the transcript of each release was fixed, no decoder became
cleverer in any way the issuing guarantees priced, and what grew was the
background corpus against which the release was decoded. The companion
theory paper's erosion result is this mechanism stated exactly, the
floor falling through H(X \textbar{} B\_t) while the leakage bound on
the release itself stands intact.

\subsubsection{5.5 Differential privacy under continual observation: the
dual
problem}\label{differential-privacy-under-continual-observation-the-dual-problem}

Continual-observation differential privacy gives privacy a time axis in
which the releases accumulate against a fixed adversary (Dwork et
al.~2010; Chan, Shi and Song 2011), with tight continual-release costs
still an active topic (Jain et al.~2023). The archival regime is the
dual: one release, fixed transcript, growing decoder. The two regimes
compose in deployed systems, but the dual problem, a fixed transcript
decoded by tomorrow's adversary, is not posed in this strand, and the
distinction is load-bearing for what ``privacy over time'' means in
each.

\subsubsection{5.6 Everlasting constructions: the t-independent
term}\label{everlasting-constructions-the-t-independent-term}

The everlasting-privacy line owns the observation that only
unconditional protection is independent of t. Moran and Naor (2006) coin
everlasting privacy for vote secrecy against later, computationally
unbounded adversaries, motivated exactly by the ageing of computational
assumptions; Aumann, Ding and Rabin (2002) achieve everlasting security
against any future advance provided the adversary's storage was bounded
at transmission time, a constructive answer to drift within its model;
Unruh (2013) gives the composable form and locates what cannot be made
everlasting classically; and Haines et al.~(2023) systematise the
e-voting protocols with everlasting privacy, the only prior SoK on
time-robust privacy guarantees we found. The long-term archiving survey
of Vigil et al.~(2015) covers renewal architectures for integrity,
authenticity and proof of existence, with confidentiality explicitly out
of scope.

Read against Definition 2, these constructions engineer terms of the
capacity sum to zero against all future classes, which is the only way
to make a term of R(t) constant in t. The research programme's related
design claim is that structural context erasure between agent
invocations, distinguished as Grade-2 forgetting (mathematically
unrecoverable) from Grade-1 hiding (recoverable with keys), removes the
archive term entirely; the Grade-1/Grade-2 naming is offered as
terminology, and the placement of an erasure term inside the R(t)
decomposition is, per the novelty fence, the only part of this not owned
by the everlasting lines. The deep version of the erasure claim, a
non-vanishing obstruction to gluing local views into a global witness,
is carried in the programme register as C86 at low confidence with its
machinery unconstructed, and is mentioned here only to state that it is
not used anywhere in this paper.

\subsubsection{5.7 The wiretap lineage and time: the missing
regime}\label{the-wiretap-lineage-and-time-the-missing-regime}

The lineage that supplies Proposition 1's machinery fixes the adversary
with the channel (Wyner 1975; Leung-Yan-Cheong and Hellman 1978; Csiszár
and Körner 1978). The closest it comes to a time axis is the
fading-channel analysis of Gopala, Lai and El Gamal (2008), where the
channel varies in time and the eavesdropper is known only statistically.
The contrast is exact and, to our knowledge, unremarked: in the fading
regime the channel varies while the adversary class is fixed; in the
archival regime the recorded channel output is fixed while the adversary
class varies. The regime in which an adversary decodes yesterday's
recording with tomorrow's decoder is outside the lineage's model, and
Definition 2 is our proposal for its missing object.

\subsubsection{5.8 Two 2026 instances (verified record, presented as
evidence)}\label{two-2026-instances-verified-record-presented-as-evidence}

Both instances below are verified public record and are presented as
evidence for the archival regime, never as contributions of this paper.

\textbf{Proof-system soundness under improved audit capability.} A
soundness flaw in the Zcash Orchard shielded protocol, a missing
constraint in variable-base scalar multiplication in halo2\_gadgets, was
present from Orchard's activation in May 2022; multiple prior audits,
including audits assisted by earlier AI models, did not find it
(BlockSec 2026; CRYPTO ISAC 2026); it was found on 2026-05-29 by Taylor
Hornby (Shielded Labs audit) using a frontier model released the
previous day (BlockSec 2026), with proof-of-concept counterfeiting
demonstrated in local testing (CRYPTO ISAC 2026); an emergency soft fork
followed on 2026-06-02 (block 3,363,426) and the NU6.2 hard fork on
2026-06-03 (block 3,364,600); no evidence of exploitation exists, and,
because of the pool's privacy properties, non-exploitation cannot be
cryptographically proven (Zcash Foundation 2026; BlockSec 2026; CRYPTO
ISAC 2026). Read against the model: the archive (the chain) was fixed;
audit capability moved with a model release; the guarantee's effective
status changed with no action by any data subject.

\textbf{Withheld capability with attested existence.} Google Quantum AI
published (2026-03-31) a roughly tenfold improvement, in spacetime
volume over best prior single-instance estimates, in Shor-algorithm
resource requirements for secp256k1, withholding circuit-level methods
and publishing instead a zero-knowledge proof of the claimed resource
counts (Babbush et al.~2026); the prior baseline series into which this
lands is public and steep (Roetteler et al.~2017; Gidney and Ekerå 2021;
Litinski 2023; Chevignard, Fouque and Schrottenloher 2025; Gidney 2025).
An independent reconstruction reached parity roughly two months later
(Schrottenloher 2026), its author noting that the original relied on a
zero-knowledge proof in place of revealed circuits; the same day, a
first-person account disclosed that the core technique had been held
internally for roughly a year under publication restriction (Gidney
2026); and an open challenge, using the published zero-knowledge
verifier as its automatic scoring filter, subsequently exceeded the
withheld benchmark (ecdsa.fail record, 2026). The metric qualification
travels with the figure wherever this episode is repeated: the tenfold
claim is the paper's own, in spacetime volume; per-axis factors reported
elsewhere differ. Read explicitly against the availability semantics of
Remark 2: on the realisable-capability reading, D\_t contained the
withheld technique from its internal discovery onward, so for the
roughly one year of restriction every public estimate of the affected
R(t) ran below its realised value and correspondingly overstated t*; the
attestation of 2026-03-31 moved the auditor's estimate, not the ratio;
and the independent reconstruction two months later is the distinct
propagation event of Remark 2, the point at which the capability became
realisable beyond its originating holder. This episode is the instance
behind Section 6.

\subsubsection{5.9 Adjacent empirics: multi-agent
leakage}\label{adjacent-empirics-multi-agent-leakage}

The measured leakage of multi-agent LLM systems is external empirical
work and enters this survey as context for why the two-channel
architecture of Section 3 is the unit of analysis: sequential
composition compounds leakage up to (2\^{}N - 1) epsilon across N
agents, with measured mutual information roughly doubling from two to
five agents (Asif and Amiri 2026), and a benchmark across 4,979 traces
on five production LLMs measures inter-agent channel leakage in the
high-sixties per-cent range of traces (El Yagoubi, Badu-Marfo and Al
Mallah 2026). These are measurements of leakage in deployed
compositions, not capacity attestations, and they bound neither term of
Definition 2; they motivate the non-collusion assumption's centrality
(Section 8).

\subsection{6 The Existence-Leak Law and the Attestation
Discount}\label{the-existence-leak-law-and-the-attestation-discount}

The second 2026 instance raises a question the disclosure literature
does not answer: what does a zero-knowledge proof of capability
existence itself disclose? The attesting paper treats withholding plus
attestation as responsible disclosure (Babbush et al.~2026) and does not
model what the existence proof leaks; the reconstruction paper reports
parity but does not quantify how much the attestation accelerated it
(Schrottenloher 2026). That gap is the territory of this section.

\subsubsection{6.1 The existence-leak
conjecture}\label{the-existence-leak-conjecture}

\textbf{Conjecture 2 (existence-leak law; programme register C81, held
open pending a second independent instance).} \emph{Candidate model.}
Fix a probability space (Ω, ℱ, P) carrying a random variable M with
values in a countable method space ℳ, representing the adversary's prior
uncertainty over which technique, if any, achieves the capability; a
cost functional c : ℳ → ℝ₊ ∪ \{∞\}, the resources the method requires
against the stated instance; and, for an attested resource level τ, the
feasibility indicator F\_τ := 1\{c(M) ≤ τ\}. A \emph{capability
attestation at level τ} is a verifiable disclosure of the event \{F\_τ =
1\} that reveals no other function of M; a zero-knowledge proof of
claimed resource counts is the instance class. \emph{Statement.} For
capability attestations, I(F\_τ; M) \textgreater{} 0: a zero-knowledge
proof that a capability is feasible leaks a nonzero upper bound, namely
τ, on the reconstruction or search difficulty of the method. \emph{Scope
fence:} the conjecture concerns capability attestations; instance
attestations, such as a proof that a particular transaction is valid,
are out of scope. \emph{Floor and object distinction:} the impossibility
of leakage-resilient zero knowledge below leakage parameter one (Garg,
Jain and Sahai 2011) shows the protocol-internal analogue cannot be
driven to zero. Formally, the two results concern different random
variables: Garg, Jain and Sahai bound the information a verifier obtains
about the witness W of a single interactive execution through a leakage
oracle, a protocol-internal quantity, whereas Conjecture 2 concerns
I(F\_τ; M), the information the public feasibility event carries about
the method across systems; neither quantity bounds the other, and the
cross-system form is, to our knowledge, unclaimed in this literature.
\emph{The named missing step.} Under the candidate model the bare
inequality is immediate whenever the prior makes F\_τ non-constant: F\_τ
is then a non-constant function of M, so I(F\_τ; M) = H(F\_τ)
\textgreater{} 0. Stated so, the law would be true and empty. The
conjectural content is operational, and the proof obligation is
correspondingly: (i) an attestation class and an adversary search model
in which the claim is non-vacuous, that is, in which conditioning on
\{F\_τ = 1\} provably and quantifiably reduces the expected search cost
of reaching parity with the attested capability; and (ii) a second
independent instance, the register's own bar for raising the
conjecture's status. \emph{Evidence:} the 2026 episode of Section 5.8,
in which parity followed attested existence in roughly two months and
the attestation's verifier became the search's scoring function.

\subsubsection{6.2 The liveness instance and deduplication side
channels}\label{the-liveness-instance-and-deduplication-side-channels}

The deduplication side-channel literature substantially anticipates the
confirmation mechanism that the programme's design surface re-derives,
and is cited here as such per the novelty fence. Harnik, Pinkas and
Shulman-Peleg (2010) showed that cross-user deduplication lets a client
confirm whether specific content already exists server-side, feasible
exactly over guessable candidate spaces. The programme's claim extends
this to content addressing: a live content address is an existence claim
about its content; under deduplication, an adversary holding a candidate
confirms existence by comparing the candidate's hash against a
known-live address; the address does not leak the content, but the
liveness of the address leaks the existence of the content, and
existence bounds the search. The scope limit is inherited unchanged from
Harnik et al.: a guessable candidate space; for high-entropy content the
upper bound still leaks while recovery may remain infeasible. The
residue not owned by the prior work, per the fence, is cross-system
corroboration and subsumption under Conjecture 2:

\textbf{Conjecture 3 (corroboration monotonicity; programme register
C93, moderate confidence).} Fix a candidate space 𝒞 for the content of a
known-live address and a corroboration model in which k independent
systems each resolve the same address; let Δ\_k(X) denote the
reconstruction difficulty of the content, valued as the adversary's
expected search cost over 𝒞 given the k resolutions. (The source of this
claim writes the difficulty as D(X); the symbol is changed here to avoid
collision with the discount functional D(a) of Conjecture 4.) The
conjecture: under a fixed corroboration model, Δ\_k(X) is monotone
non-increasing in k. \emph{Proof obligation:} the fixed corroboration
model itself: a joint law over candidate spaces and resolver behaviour
under which each additional independent resolution weakly narrows the
candidate space, together with a proof of the comparison Δ\_\{k+1\}(X) ≤
Δ\_k(X) within it. Under such a model the programme additionally expects
a steep-then-shallow profile in k, which is not asserted here beyond
monotonicity. \emph{Remark:} the genealogy-database record of Section
5.4 (Gymrek et al.~2013) is a naturally occurring instance of the
monotonicity, each new database narrowing the candidate space.

\subsubsection{6.3 The attestation discount on migration
horizons}\label{the-attestation-discount-on-migration-horizons}

\textbf{Conjecture 4 (attestation discount; programme register C84,
weaker evidentiary support than Conjectures 1 to 3).} Let 𝒜 be a space
of public feasibility attestations, preordered by specificity: a ⪯ a'
when a' verifiably discloses at least every feasibility event a
discloses, with ⊥ ∈ 𝒜 the null attestation. Call D : 𝒜 → {[}0, ∞) a
\emph{discount functional} if it satisfies (D1) normalisation, D(⊥) = 0;
(D2) monotonicity in specificity, a ⪯ a' implies D(a) ≤ D(a'); and (D3)
domination by disclosure, D(a) ≤ D(a\_full) for every a, where a\_full
is full method publication. The conjecture: Z\_b' = Z\_b − D(a): horizon
updates in Mosca-style planning are described by some discount
functional D with D(a) \textgreater{} 0 for every non-null capability
attestation, independently of any actual attack. Here Z\_b is the
planner's estimate of the collapse horizon, for which the shelf life t*
of Definition 2 is the derived referent: per Remark 2, an attestation of
an already-realised capability moves the estimate and not the derived
crossing, so the discount acts on Z\_b as an estimate; no claim is made
that attestation moves t* itself. The Mosca lineage updates its horizon
by expert belief (Mosca 2018; Mosca and Piani 2024); no found work
formalises attestation-induced discounting, and the
harvest-now-decrypt-later economics (Blanco-Romero et al.~2026) prices
collection, not attestation. \emph{Status of the axioms:} (D1) to (D3)
are what the schema requires to be well posed; they do not determine D,
and the conjecture is held at schema level until a concrete D is
identified. \emph{Proof obligation:} an elicitation or market study
identifying a discount functional consistent with observed horizon
updates causally attributable to attestations rather than to
demonstrated attacks. Stated as a conjecture, not a result, at exactly
the register's status.

\subsubsection{6.4 Limitative framings (related-work status
only)}\label{limitative-framings-related-work-status-only}

The research programme also carries two readings that frame, without
strengthening, the results above; both are register-carried framing
conjectures, recorded by their source's own fence as structural framing
rather than theorem-to-theorem reduction, and neither is used as a
premise anywhere in this paper. The first (programme register C90,
recorded as an observation with no reduction claimed) reads the
load-bearing role of the unreconstructable remainder as an inversion of
the limitative pattern in logic: an architecture whose R(t) approaches
one does not merely lose a safety margin but loses the property that
made the disclosure architecture valuable, on at least one axis. The
second (programme register C92, capped at the confidence of Conjecture
2, which it extends) reads Conjecture 2 as an undefinability phenomenon:
a disclosing system cannot confine the truth of its own feasibility
claim once that claim is corroborable elsewhere. These framings are
recorded here for completeness of the register translation. Formal-pass
decision: both are retained at related-work status only. Neither admits
a formal statement without a reduction their own source declines to
claim, no formal statement of either is made in this paper, and neither
appears in, nor is used by, the formal apparatus of Sections 3, 4 and
6.1 to 6.3.

\subsection{7 Open Problem: Countermeasure by
Divergence}\label{open-problem-countermeasure-by-divergence}

The countermeasure directions surveyed above either engineer capacity
terms to zero (Section 5.6), erase the archive term (Section 5.6), or
reduce migration time (Section 5.2). The research programme names a
fourth direction and this paper reports it strictly as an open problem,
explicitly unvalidated.

\textbf{Open problem (divergence outrunning drift; programme register
C18 to C21, held at low confidence, unmeasured).} If the data subject's
protected trajectory diverges, in the Lyapunov sense with exponent
lambda \textgreater{} 0, from any snapshot of it, then reconstruction
error from a fixed archive grows as e\^{}(lambda t), and the design goal
becomes trajectory divergence outrunning capability drift. lambda is
unmeasured; the programme records it as its most-needed number. We found
no anticipation of this direction in the swept literature, and equally
no support: the claim's confidence is low by the register's own
assessment, and it earns its place in this survey only as a precisely
stated open problem. The nearest neighbour is the agility literature's y
(Näther et al.~2024), which measures the cost of changing the
protection, not the divergence of the protected process.

\subsection{8 Threats to Validity}\label{threats-to-validity}

\textbf{Construct validity.} The capacities of Definition 2 are declared
quantities, not measurements: no standardised methodology exists for
attesting C\_S(t) or C\_M(t) against a stated adversary class, and until
one does, the deficit condition is checkable only as a declaration
audit. The multi-agent leakage measurements of Section 5.9 measure
leakage of deployed compositions, not channel capacity, and do not
operationalise the model's terms.

\textbf{Internal validity.} The evidence for Conjecture 1 is a single
year's instance cluster (Section 5.8), and for Conjecture 2 a single
episode; the register's own bar for raising Conjecture 2's status is a
second independent instance, which has not occurred. The Orchard
instance's market context is deliberately omitted: the contemporaneous
repricing is confounded by a concurrent, prominent institutional exit
and is not attributable to the flaw alone; no price figure is carried in
this paper.

\textbf{External validity.} The bound of Section 3 is stated for a
two-agent architecture under Assumptions 1 and 2; adversary classes
outside the declared class, and architectures whose channel count or
composition differs, are outside every quantitative statement in this
paper. The survey is web-only, English-only, and single-session at its
base; the disclosure-economics strand is touched only through one
analysis (CRYPTO ISAC 2026) and non-English and paywalled venues were
not systematically covered. The unification claim is made to our
knowledge, and its nearest published neighbour is read and positioned
rather than presumed: the temporal-risk model of Kagai, Branch, But and
Allen (2025), engaged in Section 5.1, is a lifetime-versus-horizon
formalisation specific to the quantum transition whose load-bearing
quantities are declared or projected inputs, so it leaves the
derived-ratio claim (N1) standing and does not reach the cross-domain
systematisation that contribution 4 claims. One bibliographically
unresolved work remains the open referee risk: a preprint reported to
formalise the same lifetime-versus-horizon comparison, not cited pending
resolution (Section 5.1); if on resolution it proves to derive rather
than declare its quantities, the derived-ratio claim narrows as well.

\textbf{Assumption failure.} If Assumption 1 fails, collusion or a third
channel recreates the inter-agent residue, the capacity-sum bound of
Proposition 1(a) is unlicensed, and the floor of Proposition 1(b) with
it: the conditional-independence residual is the entire content of the
two-channel advantage, and the multi-agent measurements of Section 5.9
show inter-agent channels leaking in a majority of measured traces when
architectures do not enforce the assumption. If Assumption 2 fails,
capacities evaluated against a weaker class than the operative adversary
understate R(t) and overstate t*. If the deficit condition is asserted
without declaration, the guarantee claim is unfounded even with both
assumptions intact, since the preconditions alone never place R below
one.

\textbf{Commitment.} A measured result against any conjecture of this
paper, including a failure of the drift of Conjecture 1 to materialise
against an instrumented archive, is filed to the research programme's
register and reported with the same prominence as a confirmation.

\subsection{9 Conclusion}\label{conclusion}

Five literatures know that privacy guarantees over fixed archives
expire; none of them says when, because none of them has an object whose
value is the expiry. This paper systematised those literatures under one
such object, the time-indexed reconstruction ratio R(t) and its derived
shelf life t*, and separated, within the guarantee, what an
architecture's preconditions buy from what only a declared and expiring
capacity deficit can buy. The apparatus is stated at its honest status:
a conditional floor proved from standard machinery with its exact
hypotheses stated and claimed only as a decomposition reading, an
adversary-ordering programme grounded at definition level whose one
structural lemma is conditional on an ordering hypothesis that remains
the programme's empirical content, four conjectures carried at their
register status with proof obligations, and one open problem whose key
number is unmeasured. The immediate next steps are the ones the statuses
dictate: a witness for Definition 3's leakage functional delivering the
joint-estimator floor of Section 4.1's residue (ii), and an instantiated
ordering, or the section's demotion; a capacity-attestation methodology
that would make the deficit condition auditable; and a second
independent instance for the existence-leak law, which is the single
observation that would most change this paper's standing.

\begin{center}\rule{0.5\linewidth}{0.5pt}\end{center}

\subsection{References}\label{references}

All entries below were resolved to a live record by the this paper
prior-art sweep (2026-07-07) or are present in the canonical
bibliography. All entries now resolve to keys in pv\_v6.bib (61 verified
entries, 0 UNVERIFIED); keys re-verified against primary records at the
A4 citation pass, 2026-07-09, with the Mosca-Piani 2025 edition added
under the L085(a) sanction.

\begin{itemize}
\tightlist
\item
  Alvim, M. S., Chatzikokolakis, K., Palamidessi, C., Smith, G.
  ``Measuring Information Leakage Using Generalized Gain Functions.''
  CSF 2012, pp.~265-279.
\item
  Alvim, M. S., Chatzikokolakis, K., McIver, A., Morgan, C.,
  Palamidessi, C., Smith, G. \emph{The Science of Quantitative
  Information Flow.} Springer, 2020.
\item
  Asif, S., Amiri, M. M. ``Information-Theoretic Privacy Control for
  Sequential Multi-Agent LLM Systems.'' arXiv:2603.05520, 2026.
\item
  Aumann, Y., Ding, Y. Z., Rabin, M. O. ``Everlasting security in the
  bounded storage model.'' IEEE Trans. Inf. Theory 48(6):1668-1680,
  2002.
\item
  Babbush, R., et al.~``Securing Elliptic Curve Cryptocurrencies against
  Quantum Vulnerabilities: Resource Estimates and Mitigations.''
  arXiv:2603.28846; ePrint 2026/625, 2026.
\item
  Barker, E. \emph{Recommendation for Key Management: Part 1 General.}
  NIST SP 800-57 Part 1 Rev.~5, 2020.
\item
  Barocas, S., Nissenbaum, H. ``Big Data's End Run around Anonymity and
  Consent.'' In \emph{Privacy, Big Data, and the Public Good}, Cambridge
  University Press, 2014.
\item
  Blanco-Romero, J., Almenares Mendoza, F., García Rubio, C., Campo, C.,
  Díaz Sánchez, D. ``On the Practical Feasibility of Harvest-Now,
  Decrypt-Later Attacks.'' arXiv:2603.01091, 2026.
\item
  BlockSec. ``Zcash Orchard Soundness Bug Analysis.'' blocksec.com,
  2026-06.
\item
  Chan, T.-H. H., Shi, E., Song, D. ``Private and Continual Release of
  Statistics.'' ACM TISSEC 14(3), art. 26, 2011.
\item
  Chevignard, C., Fouque, P.-A., Schrottenloher, A. ``Reducing the
  Number of Qubits in Quantum Factoring.'' CRYPTO 2025, Springer,
  pp.~384-415.
\item
  Cover, T. M., Thomas, J. A. \emph{Elements of Information Theory.} 2nd
  ed., Wiley-Interscience, 2006.
\item
  CRYPTO ISAC. ``The Zcash Orchard Bug and the New Economics of
  Vulnerability Discovery.'' cryptoisac.org, 2026-06.
\item
  Csiszár, I., Körner, J. ``Broadcast Channels with Confidential
  Messages.'' IEEE Trans. Inf. Theory 24(3):339-348, 1978.
\item
  Dallaire-Demers, P.-L., Doyle, W., Foo, T. ``Brace for impact: ECDLP
  challenges for quantum cryptanalysis.'' arXiv:2508.14011, 2025 (rev.
  2026).
\item
  de Montjoye, Y.-A., Hidalgo, C. A., Verleysen, M., Blondel, V. D.
  ``Unique in the Crowd: The privacy bounds of human mobility.''
  Scientific Reports 3:1376, 2013.
\item
  Dwork, C., Naor, M., Pitassi, T., Rothblum, G. N. ``Differential
  privacy under continual observation.'' STOC 2010, pp.~715-724.
\item
  ecdsa.fail challenge (Eigen Labs). Repository
  github.com/ecdsafail/ecdsafail-challenge; leaderboard, 2026.
\item
  El Yagoubi, F., Badu-Marfo, G., Al Mallah, R. ``AgentLeak: A Benchmark
  for Internal-Channel Privacy Leakage in Multi-Agent LLM Systems.''
  arXiv:2602.11510, 2026.
\item
  Erlich, Y., et al.~``Identity inference of genomic data using
  long-range familial searches.'' Science 362(6415):690-694, 2018.
  doi:10.1126/science.aau4832.
\item
  Fano, R. M. \emph{Transmission of Information: A Statistical Theory of
  Communication.} MIT Press, 1961.
\item
  Mascelli, J., Rodden, M. ``\,`Harvest Now Decrypt Later': Examining
  Post-Quantum Cryptography and the Data Privacy Risks for Distributed
  Ledger Networks.'' Federal Reserve Board, FEDS 2025-093, 2025.
  doi:10.17016/FEDS.2025.093.
\item
  Garg, S., Jain, A., Sahai, A. ``Leakage-Resilient Zero Knowledge.''
  CRYPTO 2011, LNCS 6841, pp.~297-315.
\item
  Gidney, C. ``How to factor 2048 bit RSA integers with less than a
  million noisy qubits.'' arXiv:2505.15917, 2025.
\item
  Gidney, C. ``The French have the Quantum Circuits.''
  algassert.com/post/2602, 2026-06.
\item
  Gidney, C., Ekerå, M. ``How to factor 2048 bit RSA integers in 8 hours
  using 20 million noisy qubits.'' Quantum 5:433, 2021.
\item
  Gopala, P. K., Lai, L., El Gamal, H. ``On the Secrecy Capacity of
  Fading Channels.'' IEEE Trans. Inf. Theory 54(10):4687-4698, 2008.
\item
  Gymrek, M., McGuire, A. L., Golan, D., Halperin, E., Erlich, Y.
  ``Identifying Personal Genomes by Surname Inference.'' Science
  339(6117):321-324, 2013.
\item
  Haines, T., Mosaheb, R., Müller, J., Pryvalov, I. ``SoK: Secure
  E-Voting with Everlasting Privacy.'' PoPETs 2023(1):279-293.
\item
  Harnik, D., Pinkas, B., Shulman-Peleg, A. ``Side Channels in Cloud
  Services: Deduplication in Cloud Storage.'' IEEE Security \& Privacy
  8(6):40-47, 2010.
\item
  Jain, P., Raskhodnikova, S., Sivakumar, S., Smith, A. ``The Price of
  Differential Privacy under Continual Observation.'' ICML 2023, PMLR
  v202:14654-14678.
\item
  Kagai, F., Branch, P., But, J., Allen, R. ``Harvest-Now,
  Decrypt-Later: A Temporal Cybersecurity Risk in the Quantum
  Transition.'' Telecom 6(4):100, 2025. doi:10.3390/telecom6040100.
\item
  Leung-Yan-Cheong, S., Hellman, M. ``The Gaussian Wire-Tap Channel.''
  IEEE Trans. Inf. Theory 24(4):451-456, 1978.
\item
  Litinski, D. ``How to compute a 256-bit elliptic curve private key
  with only 50 million Toffoli gates.'' arXiv:2306.08585, 2023.
\item
  Moran, T., Naor, M. ``Receipt-Free Universally-Verifiable Voting with
  Everlasting Privacy.'' CRYPTO 2006, LNCS 4117, pp.~373-392.
\item
  Mosca, M. ``Cybersecurity in an Era with Quantum Computers: Will We Be
  Ready?'' IEEE Security \& Privacy 16(5):38-41, 2018.
\item
  Mosca, M., Piani, M. \emph{Quantum Threat Timeline Report 2024.}
  Global Risk Institute / evolutionQ, 2024.
\item
  Mosca, M., Piani, M. \emph{Quantum Threat Timeline Report 2025.}
  Global Risk Institute / evolutionQ, 2025 (GRI publication page posted
  2026-03-09).
\item
  Narayanan, A., Shmatikov, V. ``Robust De-anonymization of Large Sparse
  Datasets.'' IEEE S\&P 2008, pp.~111-125.
\item
  Näther, C., Herzinger, D., Steghöfer, J.-P., Gazdag, S.-L., Hirsch,
  E., Loebenberger, D. ``Toward a Common Understanding of Cryptographic
  Agility: A Systematic Review.'' arXiv:2411.08781, 2024.
\item
  Rocher, L., Hendrickx, J. M., de Montjoye, Y.-A. ``Estimating the
  success of re-identifications in incomplete datasets using generative
  models.'' Nature Communications 10:3069, 2019.
  doi:10.1038/s41467-019-10933-3.
\item
  Roetteler, M., Naehrig, M., Svore, K. M., Lauter, K. ``Quantum
  Resource Estimates for Computing Elliptic Curve Discrete Logarithms.''
  ASIACRYPT 2017, LNCS 10625, pp.~241-270.
\item
  Schrottenloher, A. ``Optimized Point Addition Circuits for Elliptic
  Curve Discrete Logarithms.'' ePrint 2026/1128; arXiv:2606.02235, 2026.
\item
  Unruh, D. ``Everlasting Multi-Party Computation.'' CRYPTO 2013; J.
  Cryptology 31(4):965-1011, 2018.
\item
  Vigil, M. A. G., Buchmann, J., Cabarcas, D., Weinert, C., Wiesmaier,
  A. ``Integrity, authenticity, non-repudiation, and proof of existence
  for long-term archiving: A survey.'' Computers \& Security 50:16-32,
  2015.
\item
  Wyner, A. D. ``The Wire-Tap Channel.'' Bell System Technical Journal
  54(8):1355-1387, 1975.
\item
  Zcash Foundation. ``Zebra 4.5.3 and 5.0.0: Emergency Soft Fork and
  NU6.2 Activation.'' zfnd.org, 2026-06.
\end{itemize}

Not cited, pending bibliographic resolution (do not cite until resolved;
see Sections 5.1 and 8): the time-dependent threat-model preprint (no
resolvable venue or author list).

\end{document}
